% §15 Exercises
%
% Cross-reference resolution status:
%
% DEFINED here:
%   \label{ex:ch15-1} through \label{ex:ch15-11}
%
% RESOLVED (used here, defined elsewhere):
%   \Cref{def:bphf}                              -- §15.2
%   \Cref{sec:bphf-construction}                 -- §15.2
%   \Cref{alg:bphf-build-query}                  -- §15.2
%   \Cref{prop:bphf-space}                       -- §15.3
%   \Cref{sec:space-time-analysis}               -- §15.3
%   \Cref{prop:bphf-ranked-search}               -- §15.4
%   \Cref{sec:ranked-search}                     -- §15.4
%   \eqref{eq:bphf-fnr}                          -- §15.2
%   \eqref{eq:bphf-r}                            -- §15.2
%   \Cref{def:bhf}                               -- chapter 14 §14.2
%   \Cref{prop:bhf-fpr}                          -- chapter 14 §14.4
%   \Cref{thm:bhf-optimality}                    -- chapter 14 §14.3
%   \Cref{ch:bhf}                                -- chapter 14
%
% Difficulty-marker macros (matching chapters 1-14):
\providecommand{\exstar}{$\star$}
\providecommand{\exstarstar}{$\star\star$}
\providecommand{\exstarstarstar}{$\star\star\star$}

\section*{Exercises}
\addcontentsline{toc}{section}{Exercises}

\begin{enumerate}

% -----------------------------------------------------------------------
% WARM-UP  (★ × 3): exercises 1, 2, 3
% -----------------------------------------------------------------------

\item \exstar
\label{ex:ch15-1}
\textbf{BPHF space for $n = 10^6$.}
For a designated set of size $n = 10^6$, compute the rank width
$r = \lceil \log_2 n \rceil$ and the per-element bit cost the BPHF
of \Cref{def:bphf} pays for the rank table. Multiply through to
report the total rank-table cost in bits and in megabytes. Compare
the result to \Cref{prop:bphf-space} and state the order of
magnitude of the sublinear perfect-hash correction the construction
inherits from chapter~14 at this $n$.

\item \exstar
\label{ex:ch15-2}
\textbf{FNR computation at $r = 10$.}
For rank width $r = 10$ bits per slot, compute the BPHF's FNR upper
bound from \eqref{eq:bphf-fnr}. Express the result as a fraction,
as a decimal, and as a count per ten thousand member queries. State
in one sentence why \eqref{eq:bphf-fnr} is an upper bound rather
than an equality (the per-member miss probability accounts for both
rank-collision and second-stage-discrimination contributions).

\item \exstar
\label{ex:ch15-3}
\textbf{BHF vs BPHF: when to use which.}
Write one paragraph each (three to five sentences) for the two
filter families.
\begin{enumerate}
  \item[(a)] When to use the BHF of \Cref{def:bhf}. Address the
    error structure (FNR = 0, controllable FPR), the bit-budget
    formula ($\lceil -\log_2 \fprate \rceil$ bits per element from
    \Cref{thm:bhf-optimality}), and one representative application.
  \item[(b)] When to use the BPHF of \Cref{def:bphf}. Address the
    error structure (FPR = 0, controllable FNR), the bit-budget
    formula ($\lceil \log_2 n \rceil$ bits per element from
    \Cref{prop:bphf-space}), and one representative application from
    \Cref{sec:ranked-search}.
\end{enumerate}

% -----------------------------------------------------------------------
% MEDIUM  (★★ × 5): exercises 4, 5, 6, 7, 8
% -----------------------------------------------------------------------

\item \exstarstar
\label{ex:ch15-4}
\textbf{Notebook-tied BPHF build.}
Use the chapter~15 Python notebook to build a BPHF for $n = 10000$
items with rank width $r = \lceil \log_2 10000 \rceil = 14$ bits per
slot.
\begin{enumerate}
  \item[(a)] Construct a designated set $A$ of $10000$ random
    $32$-bit integers and assign each a distinct rank in $[0, n)$
    using the notebook's rank-assignment helper. Build the BPHF.
  \item[(b)] Generate $10^5$ random query points drawn from outside
    $A$ (e.g., from a disjoint integer range) and count the number
    of queries that return $\top$. Verify the count is exactly $0$,
    confirming the $\fprate = 0$ guarantee of \Cref{def:bphf}
    holds empirically.
  \item[(c)] Query each $a \in A$ once, count the misses (queries
    returning $\bot$ for known members), and compute the empirical
    FNR. Compare to the upper bound $2^{-14} \approx 6.1 \times
    10^{-5}$ from \eqref{eq:bphf-fnr}, and state whether the
    agreement is within an order of magnitude.
\end{enumerate}

\item \exstarstar
\label{ex:ch15-5}
\textbf{Top-$k$ miss rate at $k = 50, n = 10^6, r = 16$.}
For a top-$k = 50$ retrieval over a corpus of $n = 10^6$ documents
with a BPHF at $r = 16$ bits per slot (a lower-resolution
configuration than the chapter's worked example, illustrating the
trade between rank precision and miss rate), apply the bound of
\Cref{prop:bphf-ranked-search}.
\begin{enumerate}
  \item[(a)] Compute $2^{-r} = 2^{-16}$ and report the per-document
    miss probability bound.
  \item[(b)] Compute the expected number of misses per top-$k$
    query, $k \cdot 2^{-r}$.
  \item[(c)] State the operational interpretation: out of how many
    queries does the BPHF expect to miss at least one top-$k$
    document, on average? (Use the Poisson approximation: for small
    expected count $\mu$, the probability of at least one miss per
    query is approximately $\mu$.)
\end{enumerate}

\item \exstarstar
\label{ex:ch15-6}
\textbf{BPHF storage vs full sorted index.}
For $n = 10^6$ documents at rank width $r = \lceil \log_2 n \rceil =
20$ bits per slot, the BPHF rank-table cost is $n \cdot r$ bits.
\begin{enumerate}
  \item[(a)] Compute the BPHF rank-table cost in megabytes.
  \item[(b)] Compute the cost of a full sorted index that stores
    each document's rank explicitly: this is also $n \log_2 n$ bits
    for the ranks, plus addressing overhead for an associative
    structure (a hash table or balanced tree) that maps document IDs
    to ranks. State the per-document associative-structure overhead
    as roughly $32$ to $64$ bits (a pointer and a load-factor
    correction) and add it to the rank cost. Compute the total
    sorted-index cost in megabytes.
  \item[(c)] Compute the ratio of BPHF cost to sorted-index cost.
    Discuss what the BPHF buys at the same dominant-term cost: the
    $\fprate = 0$ membership check is folded into the rank lookup,
    avoiding the explicit ID-key storage that the sorted index pays
    for. The BPHF's space saving is the associative-structure
    overhead, on the order of a factor of $2$ to $4$ at $n = 10^6$.
\end{enumerate}

\item \exstarstar
\label{ex:ch15-7}
\textbf{Compose BPHF with BHF for $x \in A \cap B$.}
Can a BPHF on a set $A$ and a BHF on a set $B$ be composed to test
the membership predicate $x \in A \cap B$, that is, ``$x$ is in
$A$ and in $B$''?
\begin{enumerate}
  \item[(a)] Sketch the construction: a query for $x$ runs the BPHF
    on $A$ and the BHF on $B$ independently, and returns $\top$ if
    and only if both filters return $\top$. State the per-query
    cost as the sum of the two underlying costs ($O(1) + O(1) =
    O(1)$).
  \item[(b)] Derive the combined error rates. The composed
    filter's FPR is the probability that $x \notin A \cap B$ but
    both filters return $\top$. Case-split: if $x \in A$ and $x
    \notin B$, the FPR contribution is the BHF's FPR ($\le 2^{-r_B}$
    where $r_B$ is the BHF's width). If $x \notin A$, the BPHF
    returns $\bot$ with certainty (its FPR is $0$ from
    \Cref{def:bphf}), so this case contributes nothing. Conclude
    the composed FPR is at most $2^{-r_B}$.
  \item[(c)] Derive the composed FNR. For $x \in A \cap B$, the
    composed query returns $\top$ only if both filters return
    $\top$; the BHF returns $\top$ with certainty (FNR $= 0$), so
    the composed miss probability is the BPHF's miss probability
    on $A$, bounded by $2^{-r_A}$ from \eqref{eq:bphf-fnr}.
    Conclude the composed filter inherits BPHF FNR and BHF FPR, the
    weaker of the two on each axis.
\end{enumerate}

\item \exstarstar
\label{ex:ch15-8}
\textbf{Update cost for a dynamic ranked corpus.}
The BPHF of \Cref{def:bphf} is built once on a static $A$ with a
fixed rank assignment $\rho$. For an application where the corpus
evolves (new documents arrive, old documents are deleted, scores
get re-computed), the rank assignment changes and the BPHF must be
rebuilt or updated.
\begin{enumerate}
  \item[(a)] Identify the structural reason a single-document
    insertion forces a $\Theta(n)$ rebuild in the worst case: the
    insertion changes the rank assignment for the documents whose
    scores are below the inserted document's score, shifting all
    such ranks down by one. The BPHF's rank table must be updated
    at every affected slot, plus the perfect hash $\phi$ must be
    rebuilt to accommodate the new element.
  \item[(b)] Sketch a tiered amortization scheme: maintain a static
    BPHF on the base corpus $A_0$ and a smaller auxiliary structure
    (a hash table mapping new document IDs to their current ranks)
    on the delta $\Delta$ of recent insertions. Queries check both;
    rebuilds happen when $|\Delta|$ exceeds a threshold (say, $0.1
    n$). Estimate the amortized cost per insertion as $O(n /
    \text{threshold}) = O(10)$, plus the per-insertion auxiliary
    cost $O(1)$.
  \item[(c)] Discuss the trade-off between rebuild frequency and
    auxiliary-structure size: smaller thresholds mean smaller
    auxiliary structures (less space) but more frequent rebuilds
    (more amortized work per insertion). State the analogue of the
    chapter~14 exercise on dynamic BHF: the BPHF and BHF have the
    same dynamic-update profile because both rest on the same
    perfect-hash machinery.
\end{enumerate}

% -----------------------------------------------------------------------
% CHALLENGING / EXTENSION  (★★★ × 3): exercises 9, 10, 11
% -----------------------------------------------------------------------

\item \exstarstarstar
\label{ex:ch15-9}
\textbf{BPHF FNR exact derivation.}
The bound \eqref{eq:bphf-fnr} states $\fnrate \le 2^{-r}$ as an
upper bound. Derive the exact FNR, accounting for the rank-collision
structure the proof in \Cref{sec:bphf-construction} sketched but did
not fully unpack.
\begin{enumerate}
  \item[(a)] \emph{Rank-collision contribution.} Suppose the rank
    assignment $\rho$ is uniform over $\{0, 1, \ldots, n - 1\}$ and
    independent across members. For a fixed member $a \in A$, the
    probability that some other member $a' \in A \setminus \{a\}$
    has rank $\rho(a') = \rho(a)$ is approximately $(n - 1) \cdot
    2^{-r}$ for $n \ll 2^r$ (the birthday-type estimate). For $n
    \approx 2^r$, the probability is closer to $1 - e^{-n / 2^r}
    \approx 1 - e^{-1} \approx 0.63$.
  \item[(b)] \emph{Tie-breaking contribution.} When two members
    collide, the construction's tie-breaking rule favors one over
    the other, causing the disfavored member to miss. The
    per-collision miss contribution is $1/2$ under a fair tie-break
    or $1$ under an unfair rule (always favoring one member).
  \item[(c)] Combine the two contributions to write the exact FNR
    as a function of $n$ and $r$. Show that for $r = \lceil \log_2 n
    \rceil$ (the chapter's standard choice), the exact FNR is
    approximately $0.5 \cdot (1 - e^{-n / 2^r}) \approx 0.5 \cdot
    (1 - e^{-1}) \approx 0.32$, much larger than the headline upper
    bound $2^{-r} \approx 1/n$ would suggest, unless a second-stage
    discriminator is added. Discuss the role the second-stage
    discriminator plays in pushing the exact FNR down toward the
    headline bound.
\end{enumerate}

\item \exstarstarstar
\label{ex:ch15-10}
\textbf{Adversarial rank assignment.}
The rank assignment $\rho$ in \Cref{sec:bphf-construction} is
chosen by the application from a uniform distribution over $\{0, 1,
\ldots, n - 1\}^A$. Suppose an adversary controls $\rho$ instead.
Can the adversary force $\fnrate$ above the bound $2^{-r}$ of
\eqref{eq:bphf-fnr}?
\begin{enumerate}
  \item[(a)] \emph{Attack one.} Choose $\rho$ to maximize
    rank-collisions: assign every member the same rank $\rho(a) = 0$
    for all $a \in A$. The build pass then writes $0$ into every
    slot of $R$; the query for any $a \in A$ reads $R[\phi(a)] = 0$
    and compares to $\rho(a) = 0$, getting equality, so the query
    returns $\top$. The FNR is $0$ in this attack (every member
    passes the check), but every non-member $x$ with $\rho(x) = 0$
    also passes, breaking the FPR $= 0$ guarantee.
  \item[(b)] \emph{Attack two.} Choose $\rho$ so that some specific
    members have ranks coinciding with the chosen ranks of other
    members at the same slot, maximizing collision-induced
    tie-break misses. Estimate the FNR under this attack as a
    function of the fraction of $A$ targeted.
  \item[(c)] Conclude: adversarial $\rho$ can violate either the
    FNR bound or the FPR guarantee, depending on which the adversary
    optimizes. The chapter's analysis assumes $\rho$ is supplied by
    a non-adversarial source (an application-specified scoring
    function); the FNR bound holds under that assumption. Discuss
    what changes when $\rho$ comes from a salted hash (a
    cryptographic source) versus from an external metadata table
    (a non-adversarial but potentially correlated source).
\end{enumerate}

\item \exstarstarstar
\label{ex:ch15-11}
\textbf{Hybrid filter on the $(\fprate, \fnrate)$ plane.}
The BHF lives at $(\fprate, \fnrate) = (\fprate_0, 0)$ for some
target $\fprate_0$; the BPHF lives at $(0, \fnrate_0)$ for some
target $\fnrate_0$. Design a hybrid filter that interpolates between
the two extremes, achieving any operating point
$(\fprate, \fnrate)$ on a controlled curve in the $(\fprate,
\fnrate)$ plane.
\begin{enumerate}
  \item[(a)] \emph{Construction.} Build a BPHF on $A$ at rank width
    $r_{\text{BPHF}}$ (giving $\fnrate \le 2^{-r_{\text{BPHF}}}$
    and $\fprate = 0$), then post-process its $\top$ verdicts by
    a probabilistic filter that admits a fraction $p$ of $\top$
    verdicts to a secondary check that may admit some
    non-members. The secondary check is a BHF on a related set
    (say, $A$ extended with some non-members the application
    has labeled as ``acceptable false positives''). Adjust $p$ and
    the secondary check's parameters to slide the operating
    point.
  \item[(b)] \emph{The operating curve.} For each setting of $p$,
    the hybrid filter's $(\fprate, \fnrate)$ lies on a curve
    parameterized by $p \in [0, 1]$. At $p = 0$ the hybrid is the
    pure BPHF; at $p = 1$ the hybrid is the pure BHF (with the
    BPHF's pre-filter acting as a no-op). For intermediate $p$ the
    operating point traces a continuous curve interpolating the
    two extremes.
  \item[(c)] \emph{Cost.} Compute the bit cost of the hybrid as the
    sum of the BPHF rank table and the BHF fingerprint table, plus
    one bit per element for the routing decision (whether the
    secondary check is consulted). Compare to the cost of a pure
    BPHF and a pure BHF at the same operating point. Discuss when
    the hybrid is worth the extra bits (when the application's
    quality metric depends on a balanced combination of FPR and
    FNR rather than dominantly on one or the other).
\end{enumerate}

\end{enumerate}
